\documentclass[11pt]{article}
\usepackage{psfig}
\usepackage{latexsym}
\usepackage{graphicx}
\usepackage{amssymb}
\usepackage{mathrsfs}
\usepackage{amscd}
\usepackage{amsmath}
%\pagestyle{headings}

\usepackage{epsfig}
%\def\Real{{I\!\!R}}
%\usepackage{hanging}

\newtheorem{theorem}{Theorem}[section]
\newtheorem{corollary}[theorem]{Corollary}
\newtheorem{definition}{Definition}[section]
\newtheorem{proposition}{Proposition}[section]
\newtheorem{lemma}{Lemma}[section]
\newenvironment{proof}{{\it Proof. }}{$\triangleleft$}
\newenvironment{remark}{{\it Remark. }}{$\triangleleft$ \medskip}


\oddsidemargin 0 pt
\textheight 24cm
\voffset -1cm
\evensidemargin 0cm
\topmargin .5cm
\textheight 22cm
\textwidth 16cm
\parskip 3mm
\parindent 5mm
%\input epsf
%\input seteps


%\newtheorem{Lemma}{Lemma}
%\newtheorem{Theorem}[Lemma]{Theorem}
%\newtheorem{cor}{Corollary}
%\newtheorem{prop}[Lemma]{Proposition}
%\newtheorem{lem}[thm]{Lemma}
%\newtheorem{que}{Question}
%\newtheorem{defn}{Definition}
%\newtheorem{example}{Example}
\newcommand{\bean}{\begin{eqnarray*}}
\newcommand{\eean}{\end{eqnarray*}}
\newcommand{\bea}{\begin{eqnarray}}
\newcommand{\eea}{\end{eqnarray}}
\newcommand{\bmt}{\left[ \begin{array}{ccccccccc}}
\newcommand{\emt}{\end{array}\right]}
\newcommand{\eq}{\begin{equation}\begin{array}{lllllllllllllllllllllllllllllllll}}
\newcommand{\ee}{\end{array}\end{equation}}

\renewcommand{\theequation}{\thesection.\arabic{equation}}
% MATH -----------------------------------------------------------
\newcommand{\norm}[1]{\left\Vert#1\right\Vert}
\newcommand{\abs}[1]{\left\vert#1\right\vert}
\newcommand{\set}[1]{\left\{#1\right\}}
\newcommand{\Real}{\mathbb R}
\newcommand{\eps}{\varepsilon}
\newcommand{\To}{\longrightarrow}
\newcommand{\BX}{\mathbf{B}(X)}
\newcommand{\A}{\mathcal{A}}
% ----------------------------------------------------------------
\title{Insect Drift Movement Modeling and Parameter Estimations via Boundary Characteristics\thanks{This research was supported by NSF DEB  0515781, 1021203 and DMS 0719783.}}


\author{John D. Reeve, Xiuquan Wang, Mingqing Xiao\thanks{Corresponding author: Prof. Mingqing Xiao, Department of Mathematics, Southern Illinois University, Carbondale, IL 62901. E-mail: mxiao@siu.edu. Fax: (618) 453-5300.}, Dashun Xu\\
Southern Illinois University, Carbondale, IL 62901\\ \\
James T. Cronin\\
Louisiana State University,
Baton Rouge, LA 70803}

\begin{document}
\maketitle

% ----------------------------------------------------------------
\begin{abstract}
One of the fundamental goals of ecology is to examine how dispersal affects the distribution and dynamics of insects across natural landscapes. More often insect movements are not completely random because of surrounding environment, and  their movements are more likely  a combination of both random and biased ones. Although it is straightforward to adopt some extensions of diffusion models to represent those more complicated patterns of movement, much less is known about how to estimate the key parameters or rate constants in these models. In this paper we present some important characteristics of insect biased movement in a rectangular domain of two-dimensional space and show that if insects are released at the center of a rectangle, then the ratios of the crossing density, the flux density , and the first mean passage time of opposite boundaries are constants, which are determined by the ratios of actual drifts in horizontal and vertical directions, respectively. Different from existing approaches, these characteristics greatly simplify the estimations of the parameters (diffusion rate, advection coefficients, and death rate) and make quantification of dispersal more approachable.
\end{abstract}

{\bf Keywords:} Diffusion rate; advection coefficients; flux; crossing probability;  biased movement.

{\bf MSC2010:} Primary: 92B05, 92D25; Secondary: 92D40 92D50.
%{\bf Keywords:} Poisson equations, Five-point stencil, Matrix inverse, Infinity norm. 

\medskip
%% MSC codes here, in the form: \MSC code \sep code
%% or
%{\bf MSC2010:} 15A09, 15A39, 35J05, 39A12.





\section{Introduction}

Dispersal or movement is one of the fundamental processes affecting the distribution and dynamics of insects across natural landscapes.  The movements of many insects appear to be a combination of both random and biased ones.  For instance, when a resource (such as food plants or prey) differs significantly in quality or quantity, insects may tend to move in a preferred direction (Kareiva 1982; Kareiva and Odell 1987).  They may also bias their movements near edges or boundaries between landscape elements (Schultz and Crone 2001; Bommarco and Fagan 2002; Ovaskainen 2004; Crone and Schultz 2008; Olson and Andow 2008; Ovaskainen et al. 2008; Reeve, Cronin and Haynes 2008; Ries and Sisk 2008; Reeve and Cronin 2010).  Another common type of biased movement is the response of insects to pheromones or kairomones (Turchin and Simmons 1997; Byers and Zhang 2011; Colazza et al. 2014).  

These studies suggest models of insect movement need to incorporate both random and biased components, with these components depending on location and local conditions.  Models of insect movement, such as convection-diffusion equations, can be readily modified to incorporate these features (e.g., Kareiva 1982; Banks et al. 1988; Kareiva and Odell 1987; Bommarco and Fagan 2002; Ovaskainen et al. 2008; Reeve et al. 2008; Reeve and Cronin 2010).  A common approach to estimating the parameters in these models is to mass release insects at various locations in the landscape, and then record their positions over time.  Alternately, individual insects are followed over time as they move through the landscape.  Models are then fitted to these observations to estimate the model parameters, using likelihood or least squares.  This approach can have three potential drawbacks: (1) it is computationally intensive because it requires repeated numerical solution of the partial differential equations, (2) the insects must be located repeatedly across the landscape, and (3) the models  assume particular functional forms for the convection and diffusion terms in the model.  The form of these functions may be deduced from biological principles, or they may simply be chosen for their flexibility.  While this approach can provide an adequate fit to the observed data, it provides only indirect information on these functional forms and how they vary with location and local conditions.  

Fagan (1997) proposed a complement of existing methods for estimating the diffusion rate in a simple diffusion model.  The organisms are released within a square, and the cumulative proportion reaching the edges recorded over time.  The solution of a diffusion model predicting this quantity is then fitted to the data using nonlinear least squares, and an estimate of the diffusion rate so obtained.  This method avoids the problem of tracking individual insects or observing their density on a grid, but does require an absorbing barrier of traps completely surrounding the square.  The solution to be fitted is also complex in form.  Xiao et al. (2013) present a method of estimating the diffusion rate using the mean occupancy time for insects in a circular region.  The position of the insect is not needed for this method, only the time it leaves the circle.  This method can also be used to estimate the parameters for one type of edge behavior, a biased random walk (Ovaskainen 2004; Ovaskainen et al. 2008).  Although these methods do not require the position of the insect, and so should be less labor intensive, they only provide estimates of the diffusion rate.  They have not been adapted to scenarios where there is substantial convection or biased movement.  

In this paper, we present some important characteristics of insect movement with bias in a rectangular domain of two-dimensional space. More specifically, we show that if insects are released at the center of a rectangle, then the ratios of the crossing density (REEVE - do you mean accumulated density here?), the flux density, and the first mean passage time of opposite boundaries are constants, which are determined by the ratios of actual drift coefficients in the horizontal and vertical directions, respectively. Different from existing approaches, these characteristics greatly simplify the estimation of the model parameters (the diffusion rate, advection coefficients, and death rate (REEVE - does the paper quantify the death rate anywhere?)) and make the quantification of dispersal more convenient.  

The paper is organized as follow. In Section 2, we provide some preliminaries of formulation of the diffusion model that includes drift movement, and then the main results of this paper are presented. In Section 3, we give a detailed procedure for parameter estimation and a process of computer-based simulation to mimic the field experiment based on statistic theory is provided, including some numerical simulations. Mathematical justification of the main results are shown in Section 4. The paper ends with concluding remarks and suggestions for field studies of dispersal in Section 5. 

\section{Preliminary Formulation and Main Results}
\subsection{Preliminary Formulation}   
The simplest way to formulate the insect movement is probably via Fick's law. If we denote the flux density in a region $U$ of two- or three-dimensional space by $\vec J$, and the population density evolution in time by $u$, then the rate of change of the total quantity within any subregion $V\subset U$ equals the negative of the net flux through the boundary of $V$, i.e.,
\[\frac{d}{dt}\int_V udx=-\int_{\partial V}\vec{J}\cdot \vec{n}ds,\]
where $\vec n$ is the outward normal to the boundary of $V$. Thus we have
\bea\label{de}
u_t=-\hbox{div}\vec{J}=-\nabla\cdot \vec{J}
\eea
by the Gauss-Green Theorem as $V$ was arbitrary. If we are describing the population density $u$ that diffuse at a rate $d$ and are advected or engage in a different movement with velocity $\vec c$, then the flux can be expressed as
\[ \vec{J}=-d\nabla u +\vec{c}\;u.\]
Substituting into (\ref{de}), we have the diffusion model 
\[u_t= \nabla\cdot\left( d\nabla u\right)-\nabla\cdot(\vec{c}u).\]
If $d$ is a constant and $\vec{v}$ is a constant vector, then the two dimensional diffusion model has a form as
\bean
\frac{\partial u}{\partial t}=d\left(\frac{\partial^2u}{\partial x^2}+\frac{\partial^2u}{\partial y^2}\right)-c_h\frac{\partial u}{\partial x}-c_v\frac{\partial u}{\partial y},
\eean
where $(c_h, c_v)=\vec{c}$. 
If we further assume the mortality rate (or growth rate) given by $\gamma(t)$, then the above equation can be modified as
\bea\label{model}
\frac{\partial u}{\partial t}=d\left(\frac{\partial^2u}{\partial x^2}+\frac{\partial^2u}{\partial y^2}\right)-c_h\frac{\partial u}{\partial x}-c_v\frac{\partial u}{\partial y}+\gamma(t) u.
\eea

Because of their abundance and their high mobility as compared to animals of similar size, insects are suitable for the study of dispersal, and the study of insect dispersal plays an essential role in estimating the area spread in order to develop successful conservation strategies for threatened or endangered species, or to suppress pest populations below economically damaging levels. 
\subsection{Main Results}
Let $\Omega\subset \Real^2$ be a bounded domain and the boundary of $\Omega$ be denoted by $\partial \Omega$. Suppose that $\Gamma \subseteq \partial \Omega$ is a connected part of $\partial \Omega$ (with Lebesgue measure $m(\Gamma)>0$). Throughout this section, we make the following two assumptions:
\begin{itemize}
  \item individuals who encounter the boundary cross it immediately and thereby maintain the density on the boundary at zero;
  \item insects are released inside $\Omega$. 
\end{itemize}  
For any part $\Gamma$ of boundary, we use the following notations: 
\bean
J_\Gamma(t)&=& \int_{\Gamma}\vec{J}\cdot\vec{n}ds=\hbox{the total projection of flux density to $\vec{n}$ along $\Gamma$ at time $t$},\\
N_{\Gamma}(t)&=&\int_0^t J_{\Gamma}(\tau)d\tau=\hbox{the accumulated population  across $\Gamma$ for $0<\tau \le t$ },\\
T_{\Gamma}(t) &=&\int_0^{t} \tau\;J_\Gamma(\tau)d\tau=\hbox{the mean first passage time to $\Gamma$  for $0<\tau \le t$}.
%P_{\Gamma}(t)&=&\frac{{N}_\Gamma(t)}{N_{\partial\Omega}(t)}=\hbox{the probability of an individual across $\Gamma$ at the time $t$}.
\eean
The first two quantities are easy to understand since for those points in $\Gamma$, $\vec{J}\cdot \vec{n}$ describes  the projection of $\vec{J}$ in a direction of moving away from $\Omega$. For the third one, if we consider an insect who starts out from an initial point $P_0$ inside $\Omega$, then $J_{\Gamma}(t)$ is the probability of the insect escaping across or reaching the boundary $\Gamma$ for the first time. Thus the mathematical expectation of the first reaching time to $\Gamma$ is given by $T_{\Gamma}$. 

Let's consider a rectangular domain 
\[\Omega=\left\{(x, y)|\; a\leq x\leq b, f\leq y\leq g\;\right\},\]
and denote the side $x=a, b$ and $y=f, g$ by $A$, $B$ and $F$, $G$, respectively.  
Let insect movement be described by the diffusion model (\ref{model}). Suppose that insects are initially released at the center of $\Omega$. Then for any time instant $t>0$, we have the following ratio characteristics:
\[\frac{N_B(t)}{N_A(t)}=\frac{J_B(t)}{J_A(t)}=\frac{T_B(t)}{T_A(t)}=e^{\frac{c_h}{d}(b-a)/2},\]
and
\[\frac{N_G(t)}{N_F(t)}=\frac{J_G(t)}{J_F(t)}=\frac{T_G(t)}{T_F(t)}=e^{\frac{c_v}{d}(g-f)/2},\]
where these ratios do not depend on the mortality rate $\gamma(t)$ (or growth rate provided that the requirement of minimum patch size (REEVE - I don't understand this condition) is met), and $(b-a)/2$ and $(g-f)/2$ correspond to the released distances from the horizontal and vertical boundary, respectively. The mathematical justification will be given in section \ref{MJ}. Here we make some comments to illustrate the importance of these results:
\begin{enumerate}
  \item If $c_h=0$, we can see that
\[ \frac{N_B(t)}{N_A(t)}=\frac{J_B(t)}{J_A(t)}=\frac{T_B(t)}{T_A(t)}=1\]
which implies that the insect movement pattern on the side $A$ and the side $B$ are the same. This is consistent with our intuition since $c_h=0$ indicates that there is no drift in the horizontal direction for the insect movement. Same argument holds for the case of $c_v=0$. If $c_h\not =0$ or $c_v\not =0$, then above ratios are either exponentially increased or decreased as $c_h/d$ or $c_v/d$ increase, depending on their signs.  This implies that these ratios are all sensitive to $c_h/d$ and $c_v/d$. 
  \item All the above ratios depend on $\frac{c_h}{d}$ or $\frac{c_v}{d}$ that characterize the degree of drift in the horizontal or vertical directions. (REEVE - Not sure what you're saying here.  Do you mean that as long as the ratio $c_h/d$ is the same, you get the same values on the boundary).  For two different types of insects, for example, if
 \[ \frac{c_h^{(1)}}{d_1}=\frac{c_h^{(2)}}{d_2}, \qquad \frac{c_v^{(1)}}{d_1}=\frac{c_v^{(2)}}{d_2}\]
then the drift patterns remain the same except for one type of insects  diffuses faster than the other one and settles down in a shorter period of time if $d_1\not =d_2$.
\item Moreover, these quantity relationships make the estimations of $\frac{c_h}{d}$ or $\frac{c_v}{d}$ to be straightforward and simple. For example, if we use sticky panels on each sides of the rectangle and count the accumulated number of insects on panels for a time sequence of $t_1, t_2, \cdots, t_m$, and compute
\[\bar {n}_{BA}=\frac{1}{m}\left( \frac{N_B(t_1)}{N_A(t_1)}+\frac{N_B(t_2)}{N_A(t_2)}+\cdots+\frac{N_B(t_m)}{N_A(t_m)}\right)\]
then we have
\[e^{\frac{c_h}{d}(b-a)/2}\approx \bar{n}_{BA}\qquad\hbox{or}\quad \frac{c_h}{d}\approx\frac{2}{b-a}\ln\;\bar{n}_{BA}.\]
(REEVE - We probably need to explain that method doesn't work well with only a single time.  Is that right?)  Notice that this estimation is quite easy to be implemented without exhausting effort required to track individually the spatial positions of large number of insects demanded by most existing approaches. Therefore these ratios make quantification of dispersal more approachable.  
\item Furthermore, these ratios still hold even if the panels capture rate is not $100\%$. Suppose the devices used on each side of the panels are the same. Let us assume that the captured number of insects at time $t$ on side $A$ and $B$ are $\hat{N}_A(t)$ and $\hat{N}_B(t)$, respectively. Let $\hat{N}_{lost}(t)$ the total lost of insects (i.e., not captured by all panels) and $\hat{N}_{total}(t)$ the total number of captured insects. In terms of all captured insects, the captured rate on side $A$ is $\frac{\hat{N}_A(t)}{\hat{N}_{total}(t)}$ and on side $B$ is $\frac{\hat{N}_B(t)}{\hat{N}_{total}(t)}$. Thus the correct numbers of captured by side $A$ and side $B$ should be, respectively,
\[N_A(t)=\hat{N}_A(t)+ \hat{N}_{lost}(t)\frac{\hat{N}_A(t)}{\hat{N}_{total}(t)}, \quad N_B(t)=\hat{N}_B(t)+ \hat{N}_{lost}(t)\frac{\hat{N}_B(t)}{\hat{N}_{total}(t)}.\]
Hence, it leads to
\[\frac{N_B(t)}{N_A(t)}=\frac{\hat{N}_B(t)}{\hat{N}_A(t)}.\]
Notice here that the mortality rate $\gamma(t)$ does not affect this ratio. 
\item If we scale the size of the rectangle, say $s_1(b-a)$ and $s_2(g-f)$ where $s_1$ and $s_2$ are positive real numbers, then we have
 \[\left(\frac{N_B(t)}{N_A(t)}\right)^{s_1}=\left(\frac{J_B(t)}{J_A(t)}\right)^{s_1}=\left(\frac{T_B(t)}{T_A(t)}\right)^{s_1}=e^{\frac{c_h}{d}s_1(b-a)/2},\]
and
\[\left(\frac{N_G(t)}{N_F(t)}\right)^{s_2}=\left(\frac{J_G(t)}{J_F(t)}\right)^{s_2}=\left(\frac{T_G}{T_F}\right)^{s_2}=e^{\frac{c_v}{d}s_2(g-f)/2},\]
that will allow us to compare ratios with different sizes of rectangles.
\item We also consider the case when there is no absorbing boundary condition (free boundary), and similar ratios still hold. This implies that absorbing boundary condition does not affect the estimations of $\frac{c_h}{d}$ or $\frac{c_v}{d}$, thus these ratios are indeed  fundamental in describing the drift movement of insects. 
\end{enumerate}

We here provide plots of $N_I(t)$, $J_I(t)$, and $T_I(t)$ for all sides ($I=A, B, F$ and $G$), respectively. Form Fig. \ref{fig-N}, Fig. \ref{fig-J}, and Fig. \ref{fig-T}, one can clearly see the ratio characteristics.   

\begin{figure}[!htbp]
    \centering
    \includegraphics[width=3in]{Plot_Num_of_All_Absorbing_Boundaries_vs_T_unit_d_01_Xiao.eps}
    \includegraphics[width=3in]{Plot_Num_of_All_Absorbing_Boundaries_vs_T_unit_log_d_01_Xiao.eps}
    \caption{The plot of $(N_A(t), N_B(t), N_F(t), N_G(t))$ with $d=0.1, c_{h}=0.2, c_{v}=-0.3$ on a 1 $\times$ 1 rectangle.}
    \label{fig-N}
\end{figure}

\begin{figure}[!htbp]
    \centering
    \includegraphics[width=3in]{Plot_Flux_of_All_Absorbing_Boundaries_vs_T_unit_d_01_Xiao.eps}
    \includegraphics[width=3in]{Plot_Flux_of_All_Absorbing_Boundaries_vs_T_unit_log_d_01_Xiao.eps}
    \caption{The plot of $(J_A(t), J_B(t), J_F(t), J_G(t))$ with $d=0.1, c_{h}=0.2, c_{v}=-0.3, N_{0}=1$ on a 1 $\times$ 1 rectangle to show the similarities.}
    \label{fig-J}
\end{figure}


\begin{figure}[!htbp]
    \centering
    \includegraphics[width=3in]{Plot_MOT_of_All_Absorbing_Boundaries_vs_T_unit_d_01_Xiao.eps}
    \includegraphics[width=3in]{Plot_MOT_of_All_Absorbing_Boundaries_vs_T_unit_log_d_01_Xiao.eps}
    \caption{The plot of $(T_A(t), T_B(t), T_F(t), T_G(t))$ with $d=0.1, c_{h}=0.2, c_{v}=-0.3, N_{0}=1$ on a 1 $\times$ 1 rectangle to show the similarities.}
    \label{fig-T}
\end{figure}

\section{Parameter Estimations and Statistical Experiments} 
\subsection{Absorbing boundary condition}
In this section, we consider insect movement in a given rectangular area 
\[\Omega=\left\{(x, y)|\; a\leq x\leq b, f\leq y\leq g\;\right\}\]
with absorbing boundary condition. We denote $u(P; t|P_0)$ the population density function (at time $t$) of an insect who starts out from an initial point $P_0$ inside $\Omega$. Without loss of generality, we assume $\gamma(t)\equiv 0$. Then according to lemma \ref{lem1}, $u$ can be expressed as
{\small
\eq\label{sol1}
u(P, t|P_0)=\displaystyle\frac{1}{|\Omega|}e^{\frac{c_h(x-x_0)+c_v(y-y_0)}{2d}}\sum_{n, m=1}^{\infty}
C_{(n, m)}\sin\left[\frac{n\pi(x-a)}{b-a}\right]\sin\left[\frac{m\pi(y-f)}{g-f}\right]e^{-\Lambda_{(n, m)}t}
\ee}
where  $|\Omega|$ stands for the area of $\Omega$, i.e., $|\Omega|=(b-a)(g-f)$, and
{\small
\bean
C_{(n, m)}=\sin\left[\frac{n\pi(x_0-a)}{b-a}\right]\sin\left[\frac{m\pi(y_0-f)}{g-f}\right],
\eean}
and
{\small
\bean
\Lambda_{(n, m)}=d\pi^2\left[\frac{n^2}{(b-a)^2}+\frac{m^2}{(g-f)^2}\right]+\frac{c_h^2+c_v^2}{4d}.
\eean}
Thus the population densities across each side of $\Omega$ can be computed by (\ref{sol1}). Let the released point $P_0$ be at the center of $\Omega$,  at any observation instance  time $t>0$ we know
\bea\label{es1}
\frac{N_B(t)}{N_A(t)}=e^{\frac{c_h}{2d}(b-a)}\qquad\hbox{and}\qquad \frac{N_G(t)}{N_F(t)}=e^{\frac{c_v}{2d}(g-f)},
\eea
which give the values $\frac{c_h}{2d}$ and $\frac{c_v}{2d}$, respectively. Then we can use one of the expressions of $N_A(t), N_B(t), N_F(t)$ or $N_G(t)$ to solve for $d$. 


Next we examine the statistical properties of these new estimators. Let $t_1, t_2, \dots, t_m$ be a sequence of observation time instance. Then the probability for insects to cross the boundary $A$ (up to the observation time $t$) is given by
\bean
P_A(t)=\frac{N_A(t)}{N_A(t)+N_B(t)+N_F(t)+N_G(t)},
\eean
and the other sides can be defined similarly. It is easy to check that
\bean
\frac{d}{dt}P_A(t)=\frac{d}{dt}P_B(t)=\frac{d}{dt}P_F(t)=\frac{d}{dt}P_G(t)=0\qquad \hbox{for}\; t>0.
\eean 
Thus we denote $P_A(t_i)=P_A, P_B(t_i)=P_B, P_F(t_i)=P_F$, and $P_G(t_i)=P_G$. For a given set of $(d, c_h, c_v)$, we are always to be able to generate the crossing probabilities $P_A$, $P_B$, $P_F$, and $P_G$ on each side, respectively.  

Now we assume $(P_A, P_B, P_F, P_G)$ is known. We are next going to estimate the parameters $(d, c_h, c_v)$ based on these probabilities by our estimator. First, we use $P_A$, $P_B$, $P_F$, $P_G$ to generate the multinominal random (observation) numbers at $t=t_i$ to each side of $\Omega$ provided that the initial number of insects is given, denoted by $N_0$. Then we can generate an observational data set (accumulated number of insects at $t=t_i$) 
\[\left\{\left(\hat{N}_A(t_i),\; \hat{N}_B(t_i),\; \hat{N}_F(t_i),\; \hat{N}_G(t_i)\right):\; i=1,2, \dots m\right\},\]
which represents observation samples (observed at $t=t_1, \dots, t_m$, respectively). Now we calculate the means of the ratios of opposite sides, that is,
\[ \bar{n}_{BA}=\frac{1}{m}\sum_{i=1}^m\frac{\hat{N}_B(t_i)}{\hat{N}_A(t_i)}, \quad\hbox{and}\quad \bar{n}_{GF}=\frac{1}{m}\sum_{i=1}^m\frac{\hat{N}_G(t_i)}{\hat{N}_F(t_i)}.\]
By taking advantage of the time invariance (\ref{es1}), we set
\[\bar{n}_{BA}=e^{\frac{\hat{c}_h}{2\hat{d}}(b-a)} \quad\hbox{and}\quad  \bar{n}_{GF}=e^{\frac{\hat{c}_v}{2\hat{d}}(g-f)}\]
to obtain the estimations $\bar{m}_h=\frac{\hat{c}_h}{\hat{2d}}$ and $\bar{m}_v=\frac{\hat{c}_v}{\hat{2d}}$, respectively. Next at each time $t=t_i$, By making use of the insect numbers on each side, we solve for $\hat{d}_I(i), I=A, B, F, G$. For example, given $\hat{N}_A(t_i)$ on side A, one can solve for $\hat {d}_A(i)$ by the following equation:
\bean
\hat{N}_A(t_i)/N_0= \frac{4}{(b-a)^2} e^{\frac{\bar{m}_h(a-b)}{2}}\sum_{n, m=1}^{\infty}C_{(n, m)}\frac{n\left(e^{\frac{\bar{m}_v(f-g)}{2}}-(-1)^me^{\frac{\bar{m}_v(g-f)}{2}}\right)}{m\left[1+\left(\bar{m}_v\frac{(g-f)}{m\pi} \right)^2\right] (\Lambda_{(n,m)}/\hat{d})}\left(1-e^{-\hat{d}(\Lambda_{(n, m)}/\hat{d})t_i}\right),
\eean
here 
\bean
\Lambda_{(n,m)}/\hat{d}=\pi^2\left[\frac{n^2}{(b-a)^2}+\frac{m^2}{(g-f)^2}\right]+\frac{\hat{c}_h^2+\hat{c}_v^2}{4\hat{d}^2}=\pi^2\left[\frac{n^2}{(b-a)^2}+\frac{m^2}{(g-f)^2}\right]+\bar{m}_h^2+\bar{m}_v^2
\eean
is known. For the other sides, the expressions are, respectively, 
\bean
\hat{N}_B(t_i)/N_0= \frac{4}{(b-a)^2} e^{\frac{\bar{m}_h(b-a)}{2}}\sum_{n, m=1}^{\infty}C_{(n, m)}\frac{n\left(e^{\frac{\bar{m}_v(f-g)}{2}}-(-1)^me^{\frac{\bar{m}_v(g-f)}{2}}\right)}{m\left[1+\left(\bar{m}_v\frac{(g-f)}{m\pi} \right)^2\right] (\Lambda_{(n,m)}/\hat{d})}\left(1-e^{-\hat{d}(\Lambda_{(n, m)}/\hat{d})t_i}\right),
\eean
and
\bean
\hat{N}_F(t_i)/N_0= \frac{4}{(g-f)^2} e^{\frac{\bar{m}_v(f-g)}{2}}\sum_{n, m=1}^{\infty}C_{(n, m)}\frac{m\left(e^{\frac{\bar{m}_h(a-b)}{2}}-(-1)^me^{\frac{\bar{m}_h(b-a)}{2}}\right)}{n\left[1+\left(\bar{m}_h\frac{(b-a)}{m\pi} \right)^2\right] (\Lambda_{(n,m)}/\hat{d})}\left(1-e^{-\hat{d}(\Lambda_{(n, m)}/\hat{d})t_i}\right),
\eean
and
\bean
\hat{N}_G(t_i)/N_0= \frac{4}{(g-f)^2} e^{\frac{\bar{m}_v(g-f)}{2}}\sum_{n, m=1}^{\infty}C_{(n, m)}\frac{m\left(e^{\frac{\bar{m}_h(a-b)}{2}}-(-1)^me^{\frac{\bar{m}_h(b-a)}{2}}\right)}{n\left[1+\left(\bar{m}_h\frac{(b-a)}{m\pi} \right)^2\right] (\Lambda_{(n,m)}/\hat{d})}\left(1-e^{-\hat{d}(\Lambda_{(n, m)}/\hat{d})t_i}\right).
\eean
Then we set
\[\hat{d}(i)=\frac{\hat{d}_A(i)+\hat{d}_B(i)+\hat{d}_F(i)+\hat{d}_G(i)}{4}.\]
Now for this $\hat{d}(i)$, by using
\[
\bar{n}_{BA}=e^{\frac{\hat{c}_h}{2\hat{d}(i)}(b-a)} \quad\hbox{and}\quad  \bar{n}_{GF}=e^{\frac{\hat{c}_v}{2\hat{d}(i)}(g-f)}\]
we obtain an estimate of $(\hat{d}(i), \hat{c}_h(i), \hat{c}_v(i))$ for $t=t_i$. Finally, we set
\bean
\hat{d}=\frac{1}{m} \sum_{i=1}^m \hat{d}(i), \quad\hat{c}_h=\frac{1}{m} \sum_{i=1}^m \hat{c}_h(i),\quad \hat{c}_v=\frac{1}{m} \sum_{i=1}^m \hat{c}_v(i).
\eean
The above process  provides a computer-based approach to mimic the field experiments.

Let's illustrate this process through numerical examples. Suppose that the diffusion rate $d=0.1$, and the advection coefficients $c_h=0.2$, $c_v=-0.3$, respectively, and the observation time sequence is $T=0.2, 0.4, 0.6, 0.8, 1$. Then from
\[N_I(t_i)=\int_0^{t_i}J_I(t)dt\qquad \hbox{where}\; I=A, B, F, \hbox{or}\; G\]
the crossing probabilities on each side can be calculated (see Table 1).


By using above crossing probabilities, let the initial number of released insect $N_0=100$ one can generate random numbers of insect crossing at each observation time instance $t=t_i$, which is shown in Table 1. 
\begin{table}[!htbp]
\begin{center} % used for centering table
\begin{tabular}{ccccc} % centered columns (4 columns)
\hline\hline                        %inserts double horizontal lines
 Crossing probability    & $P_A$ & $P_B$ & $P_F$ & $P_G$\\
\hline
$t=t_i$ & 0.1269  &  0.3450  &  0.4318  &  0.0963   \\ 
\hline\hline                        %inserts double horizontal lines
 Number of observed insects  & $\hat{N}_A(t_i)$ & $\hat{B}(t_i)$ & $\hat{F}(t_i)$ & $\hat{G}(t_i)$\\
\hline           
 $t_1=0.2$ & 1   &  1  &   2  &   0  \\
 $t_2=0.4$ & 3   & 12  &  12  &   4\\
 $t_3=0.6$ & 6   & 21  &  20  &   8 \\
 $t_4=0.8$ & 8   & 28  &  26  &   9\\
 $t_5=1.0$ & 10  &  32 &   30 &    9\\
\hline
\end{tabular}
\caption{The number of insects observed on each side of the boundary. The initial number of released insects is $N_0=100$. }
\label{table:nonlin} % is used to refer this table in the text
\end{center}
\end{table}
Then we set
\eq\label{rae}
\bar{n}_{BA}=e^{\frac{\hat{c}_h}{2\hat{d}(i)}(b-a)} \quad\hbox{and}\quad  \bar{n}_{GF}=e^{\frac{\hat{c}_v}{2\hat{d}(i)}(g-f)}.
\ee
Next we apply $\bar{n}_{BA}$ and $\bar{n}_{GF}$ and $t_i$ to solve $\hat{d}_I(i), I=A, B, F, G$ by setting $\hat{N}_I(i)$ as a function of $\hat{d}_I$, and compute the mean of $\hat{d}_I(i)$ as 
\[\hat{d}(i)=\frac{\hat{d}_A(i)+\hat{d}_B(i)+\hat{d}_F(i)+\hat{d}_G(i)}{4}.\]
Now combining with (\ref{rae}), we obtain an estimate $(\hat{d}(i), \hat{c}_h(i), \hat{c}_v(i))$ given in Table 2.
\begin{table}[!htbp]
\begin{center} % used for centering table
\begin{tabular}{cccc} % centered columns (4 columns)
\hline\hline                        %inserts double horizontal lines
 Parameters & $\hat{d}(i)$ & $\hat{c}_{h}(i)$ & $\hat{c}_{v}(i)$ \\
\hline
 $t_1=0.2$  &  0.0967  &  0.2151  & -0.2491 \\
 $t_2=0.4$  &  0.1008  &  0.2241  & -0.2596\\
 $t_3=0.6$  &  0.1076  &  0.2393  & -0.2772 \\
 $t_4=0.8$  &  0.1061  &  0.2360  & -0.2733 \\
 $t_5=1.0$  &  0.1039  &  0.2310  & -0.2675  \\  
\hline\hline %inserts single line\\
Mean & 0.1030& 0.2291&-0.2653 \\
\end{tabular}
\caption{The estimations of $d, c_{h}$ and $c_{v}$ and $N_0=100$. }
\label{table:nonlin} % is used to refer this table in the text
\end{center}
\end{table}


We repeat this experiment for $N_0=300, 500, 1000$, the estimation of $(d, c_h, c_v)$ as well as the standard deviation and the $95\%$ confidence intervals are provided in Table 3. REEVE - So, there is one replicate with 100 insects released, then 300, 500, and 1000?
\begin{table}[!htbp]
\centering % centering table
\begin{tabular}{ccccc} % creating 5 columns
\hline\hline % inserting double-line
Parameter &Samples & Mean & SD & 95$\%$ CI
\\ [0.5ex]
\hline % inserts single-line
% Entering 1st row
              &100  & 0.1030 & 0.0044 & (0.0967, 0.1076) \\
              &300  & 0.0975 & 0.0027 & (0.0942, 0.1008) \\
$\hat{d}$     &500  & 0.0983 & 0.0029 & (0.0953, 0.1028)  \\
              &1000 & 0.1007 & 0.0007 & (0.0998, 0.1017)  \\
\hline
              &100  & 0.2291 & 0.0097 & (0.2151, 0.2393) \\
              &300  & 0.2446 & 0.0069 & (0.2364, 0.2529) \\
$\hat{c}_{h}$ &500  & 0.2081 &0.0061& (0.2016, 0.2174)   \\
              &1000 & 0.2156 &0.0015& (0.2137, 0.2177)   \\
\hline
              &100  & -0.2653 & 0.0113 & (-0.2772, -0.2491) \\
              &300  & -0.3105 & 0.0087 & (-0.3210, -0.3001) \\
$\hat{c}_{v}$ & 500 & -0.3356 &0.0099& (-0.3507, -0.3251) \\
              &1000 & -0.2916 &0.0021& (-0.2945, -0.2891) \\
% [1ex] adds vertical space
\hline\hline % inserts single-line
\end{tabular}
\caption{Mean, standard deviation and 95$\%$ confidence interval of $(\hat{d}, \hat{c}_{h}, \hat{c}_{v})$ for absorbing boundary case.}
\label{tab:Estimation_of_Parameters_D_ch_cv}
\end{table}

\subsection{Free boundary condition}
Next we conduct similar statistical experiment for the biased movement in a rectangle with free boundary condition. We consider insect movement describing by the following diffusion model 
\bean
\frac{\partial u}{\partial t}=d\left(\frac{\partial^2u}{\partial x^2}+\frac{\partial^2u}{\partial y^2}\right)-c_h\frac{\partial u}{\partial x}-c_v\frac{\partial u}{\partial y}, \quad (x, y)\in \Real^2,
\eean
with initial condition $u|_{t=0}=\delta(P-P_0)$. In this case, the density function is given by
\bean
u(P, t|P_0)=\frac{1}{4\pi d t}\;e^{-\frac{(x-x_0-c_ht)^2+(y-y_0-c_vt)^2}{4dt}},\quad t>0
\eean
where $P_0=(x_0, y_0)$ is the released point. 
% Let $\Omega\Real^2$ be any boundary domain. For any part $\Gamma$ of the boundary of $\Omega$, the following notations are adopted
%\bean
%u(P, t|P_0)\Big|_\Gamma&=&\hbox{the population density at $P\in\Gamma$ at time $t$},\\
%{\cal N}_\Gamma(t) &=& \int_\Gamma u(P, t|P_0)\Big|_\Gamma d\Gamma=\hbox{the total population on $\Gamma$ at time $t$},\\
%{N}_\Gamma(t) &=& \int_0^t {\cal N}_\Gamma(\tau)d\tau =\hbox{the accumulated insect population on $\Gamma$ up to $t$},\\
%&=& \hbox{the insect mean occupancy time on $\Gamma$ (up to $t$)}:={\cal T}_\Gamma(t).
%\eean
Let $\Omega=[0, 1]\times [0, 1]$ and insects be released at the center of $\Omega$, i.e., $P_0=(1/2, 1/2)$. We denote the (accumulated) population  up to time $t$ on each side by ${N}_\Gamma(t)$
\[N_\Gamma(t)=\int_0^t\left(\int_\Gamma u(P, t|P_0) d\Gamma\right) dt\]
 where $\Gamma=A, B, F$, or $G$. It is straightforward to show (given in next section) that the following ratios hold:
\bean
\frac{N_B(t)}{N_A(t)}=e^{\frac{c_h}{d}(b-a)/2}, \qquad \frac{{N}_G(t)}{{N}_F(t)}=e^{\frac{c_h}{d}(g-f)/2}.
\eean
REEVE - To do this in the field, we would need a trap or device that integrates insect density over time for a defined area.  Is that right?  

Again let $d=0.1, c_h=0.2, c_v=-0.3$, and the observational time sequence is $t=0.2, 0.4, 0.6, 0.8. 1.0$. Then the probability for (accumulated) insects to be on each side at time $t=t_i$ can be calculated. 
Thus one can always generate multinomial random observations at $t=t_i$. 

For a given set of random observations:
\bean
\{\hat{{ N}}_A(t_i), \hat{{ N}}_B(t_i), \hat{{N}}_F(t_i), \hat{{ N}}_G(t_i): 1\le i\le m\},
\eean
similar to previous approach, we calculate the means of the ratios on opposite sides as following:
\[{\bar n}_{BA}=\frac{1}{m}\sum_{i=1}^m \frac{\hat{{ N}}_B(t_i)}{\hat{{ N}}_A(t_i)}, \quad {\bar n}_{GF}=\frac{1}{m}\sum_{i=1}^m \frac{\hat{{ N}}_G(t_i)}{\hat{{ N}}_F(t_i)},\]
and set
\[{\bar n}_{BA}=e^{\frac{c_h}{2d}}, \quad {\bar n}_{GF}=e^{\frac{c_h}{2d}}\]
to obtain the estimates $\bar{m}_h={\hat{c}_h}/{2\hat{d}}$ and $\bar{m}_v={\hat{c}_v}/2{\hat{d}}$. 

Next at each time instance $t=t_i$ we solve for $\hat{d}$ on each side. Notice that 
\bean
-\frac{(x-x_0-c_ht_i)^2+(y-y_0-c_vt_i)^2}{4dt_i}=-\frac{(x-x_0)^2+(y-y_0)^2}{4dt_i}+\frac{c_h(x-x_0)+c_v(y-y_0)}{2d}-\frac{(c_h^2+c_v^2)t_i}{4d}
\eean
  which is approximated by
\bean
K(x, y)=-\frac{(x-x_0)^2+(y-y_0)^2}{4dt_i}+\bar{m}_h(x-x_0)+\bar{m}_v(y-y_0)-d(\bar{m}_h^2+\bar{m}_v^2)t_i.
\eean
Now we solve for $\hat{d}_A(i)$, $\hat{d}_B(i)$, $\hat{d}_F(i)$, and $\hat{d}_G(i)$ by letting
\bean
{\hat {\cal N}}_A(t_i)/N_0=\int_0^1 e^{K(0, y)}dy, \qquad \hat{{\cal N}}_B(t_i)/N_0=\int_0^1 e^{K(1, y)}dy,\\
\hat{{\cal N}}_F(t_i)/N_0=\int_0^1 e^{K(x, 0)}dx, \qquad \hat{{\cal N}}_G(t_i)/N_0=\int_0^1 e^{K(x, 1)}dx.
\eean
Then as before,  we set $\hat{d}(i)$ to be
\[\hat{d}(i)=\frac{\hat{d}_A(i)+\hat{d}_B(i)+\hat{d}_F(i)+\hat{d}_G(i)}{4}.\]
Finally we apply the relationship 
\bean
\frac{\hat{c}_h(i)}{2\hat{d}(i)}=\bar{m}_h, \quad \frac{\hat{c}_v(i)}{2\hat{d}(i)}=\bar{m}_v
\eean
to get $\hat{c}_h(i)$ and $\hat{c}_v(i)$. Then the estimations of $(d, c_h, c_v)$ are given by
\bean
\hat{d}=\frac{1}{m}\sum_{i=1}^m \hat{d}(i), \quad \hat{c}_h=\frac{1}{m}\sum_{i=1}^m \hat{c}_h(i), \quad \hat{c}_v=\frac{1}{m}\sum_{i=1}^m \hat{c}_v(i). 
\eean 
For the free boundary case, insect movement (related to the boundary) appears to be more random and small size of insect numbers seems to be inadequate (numerical evidence also confirms that). Thus we use sample sizes $N_0=300, 500$ and $1000$.
\begin{table}[!htbp]
\centering % centering table
\begin{tabular}{ccccc} % creating 5 columns
\hline\hline % inserting double-line
Parameter &Samples & Mean & SD & 95$\%$ CI
\\ [0.5ex]
\hline % inserts single-line
% Entering 1st row
              &300  & 0.1015 & 0.0065 & (0.0950, 0.1096)  \\
$\hat{d}$     &500  & 0.0955 & 0.0041 & (0.0912, 0.1019)  \\
              &1000  & 0.0999 & 0.0017 & (0.0979, 0.1022)  \\
\hline
              &300  & 0.1876 & 0.0119 & (0.1756, 0.2026)  \\
$\hat{c}_{h}$ &500  & 0.1824 &0.0077& (0.1742, 0.1948)   \\
              &1000  & 0.2143 & 0.0037 & (0.2100, 0.2191)  \\
\hline
              &300  & -0.2706 & 0.0172 & (-0.2922, -0.2533)  \\
$\hat{c}_{v}$ &500 & -0.2919 &0.0124& (-0.3116, -0.2788) \\
              &1000  & -0.2964 & 0.0051 & (-0.3030, -0.2905)  \\
% [1ex] adds vertical space
\hline\hline % inserts single-line
\end{tabular}
\caption{Mean, standard deviation and 95$\%$ confidence interval of $\hat{d}, \hat{c_{h}}$ and $\hat{c_{v}}$ for free boundary case.}
\label{tab:Estimation_of_Parameters_D_ch_cv}
\end{table}


\section{Mathematical Justification}\label{MJ}
In this section we provide detailed mathematical justification to support our statement of previous sections. We consider the following two-dimensional diffusion model
\bea\label{two-D-model}
\frac{\partial u}{\partial t}=d\left(\frac{\partial^2u}{\partial x^2}+\frac{\partial^2u}{\partial y^2}\right)-c_h\frac{\partial u}{\partial x}-c_v\frac{\partial u}{\partial y}+\gamma(t) u, \quad (x, y)\in \Omega, \;
u|_{\partial \Omega}=0,
\eea
where $\Omega=\left\{(x, y)|\; a\leq x\leq b, f\leq y\leq g\;\right\}$.
\begin{lemma}\label{lem1} Let the initial condition be given by $u|_{t=0}=\delta(P-P_0)$, where $P=(x, y), P_0=(x_0, y_0)$ inside $\Omega$. Then the solution of (\ref{two-D-model}) is given by
{\small
\eq\label{sol}
u(P, t|P_0)=\displaystyle\frac{1}{|\Omega|}e^{\frac{c_h(x-x_0)+c_v(y-y_0)}{2d}}\sum_{n, m=1}^{\infty}
C_{(n, m)}\sin\left[\frac{n\pi(x-a)}{b-a}\right]\sin\left[\frac{m\pi(y-f)}{g-f}\right]e^{-\Lambda_{(n, m)}t+\int_0^t\gamma(\tau)d\tau}
\ee}
where  $|\Omega|$ stands for the area of $\Omega$, i.e., $|\Omega|=(b-a)(g-f)$, and
{\small
\bean
C_{(n, m)}=\sin\left[\frac{n\pi(x_0-a)}{b-a}\right]\sin\left[\frac{m\pi(y_0-f)}{g-f}\right],
\eean}
and
{\small
\bean
\Lambda_{(n, m)}=d\pi^2\left[\frac{n^2}{(b-a)^2}+\frac{m^2}{(g-f)^2}\right]+\frac{c_h^2+c_v^2}{4d}.
\eean}
\end{lemma}
\begin{proof} It is easy to see that a general solution of (\ref{two-D-model}) can be expressed as
\bean
u(x, y; t)=v(x, y; t)e^{\frac{c_h(x-x_0)+c_v(y-y_0)}{2d}-\frac{c_h^2+c_v^2}{4d}t+\int_0^t\gamma(\tau)d\tau}
\eean 
where $v$ satisfies 
\[ \frac{\partial v}{\partial t}=d\left(\frac{\partial^2v}{\partial x^2}+\frac{\partial^2v}{\partial y^2}\right), \quad (x, y)\in \Omega, \;
v|_{\partial \Omega}=0.\]
Thus one can use the method of separation of variables to obtain a general solution of $v$ as
\[v(x, y; t)= \frac{1}{|\Omega|}\sum_{n,m =1}^{\infty}\hat{C}_{(n, m)}\sin\left[\frac{n\pi(x-a)}{b-a}\right]\sin\left[\frac{m\pi(y-f)}{g-f}\right] e^{-\lambda_{(n, m)}t}\]
where 
\[ \lambda_{(n, m)}= d\pi^2\left[\frac{n^2}{(b-a)^2}+\frac{m^2}{(g-f)^2} \right].\]
Notice that the initial condition for $v$ can be expressed as
\[v(x, y; 0)=\delta(x-x_0, y-y_0)e^{-\frac{c_h(x-x_0)+c_v(y-y_0)}{2d}}\]
which gives 
\bean
\hat{C}_{(n, m)}=\sin\left[\frac{n\pi(x_0-a)}{b-a}\right]\sin\left[\frac{m\pi(y_0-f)}{g-f}\right]=C_{(n, m)}
\eean
Assembling all terms yields the expression (\ref{sol}) of $u$. The proof of uniform convergence of (\ref{sol}) for $t>0$ is standard (see, e.g., Pinsky 1998) and we thus skip it.    
\end{proof}

Notice that if $\gamma(t)$ is the mortality rate, we have $\gamma(t)\leq 0$. Thus the system is always stable. If $\gamma(t)$ is the growth rate, then we assume that the size of $\Omega$ is chosen so that
\[  \Lambda_{(1, 1)}-\sup_{t>0}\frac{1}{t}\int_0^t\gamma(\tau)d\tau >0,\]
that is, a minimum patch size is required to support a population in (\ref{two-D-model}) for this case. 

\begin{theorem}\label{thm1} Let $N_{\Gamma}$, $J_\Gamma$, $T_\Gamma$, and $P_\Gamma$ be defined as the beginning of previous section and $\Omega=\left\{(x, y)|\; a\leq x\leq b, f\leq y\leq g\;\right\}$. We denote the side $x=a, x=b$ and $y=f, y=g$ of $\Omega$ by $A$, $B$ and $F$, $G$, respectively. If $x_0=(a+b)/2$ and $y_0=(f+g)/2$, then for any  $t>0$, we have
{\small
\bea\label{ch1}
\frac{N_B(t)}{N_A(t)}=\frac{J_B(t)}{J_A(t)}=\frac{T_B(t)}{T_A(t)}=e^{\frac{c_h}{2d}(b-a)},
\eea}
and
{\small
\bea\label{ch2}
\frac{N_G(t)}{N_F(t)}=\frac{J_G(t)}{J_F(t)}=\frac{T_G(t)}{T_F(t)}=e^{\frac{c_v}{2d}(g-f)}.
\eea}
\end{theorem}
\begin{proof} We will only show (\ref{ch1}) and the proof of (\ref{ch2}) can be carried out in the same way. According to (\ref{sol}) in Lemma \ref{lem1}, a direct calculation yields
\bean
J_A(t)&=&\int_A\vec{J}\cdot\vec{n}ds=\int_Ad\frac{\partial u}{\partial x}\Big|_{x=a} ds=d\int_f^g\frac{\partial u}{\partial x}\Big|_{x=a}dy\\
&=& \frac{4d}{(b-a)^2} e^{\frac{c_h(a-x_0)}{2d}}\sum_{n, m=1}^{\infty}C_{(n, m)}\frac{n\left(e^{\frac{c_v(f-y_0)}{2d}}-(-1)^me^{\frac{c_v(g-y_0)}{2d}}\right)}{m\left[1+\left(\frac{c_v(g-f)}{2md\pi} \right)^2\right]}e^{-\Lambda_{(n, m)}t+\int_0^t\gamma(\tau)d\tau},
\eean
where $C_{(n, m)}$ and $\Lambda_{(n, m)}$ are given by Lemma \ref{lem1}. Similarly, one can have
\bean
J_B(t)&=&\int_B\vec{J}\cdot\vec{n}ds=\int_Bd\frac{\partial u}{\partial x}\Big|_{x=b} ds=-d\int_f^g\frac{\partial u}{\partial x}\Big|_{x=b}dy\\
&=& \frac{4d}{(b-a)^2} e^{\frac{c_h(b-x_0)}{2d}}\sum_{n, m=1}^{\infty}C_{(n, m)}(-1)^{n+1}\frac{n\left(e^{\frac{c_v(f-y_0)}{2d}}-(-1)^me^{\frac{c_v(g-y_0)}{2d}}\right)}{m\left[1+\left(\frac{c_v(g-f)}{2md\pi} \right)^2\right]}e^{-\Lambda_{(n, m)}t+\int_0^t\gamma(\tau)d\tau}.
\eean
Notice that the differences between $J_A(t)$ and $J_B(t)$ are the beginning terms $e^{\frac{c_h(a-x_0)}{2d}}$, $e^{\frac{c_h(b-x_0)}{2d}}$ and the additional term $(-1)^{n+1}$ appeared inside the series of $J_B(t)$. If $x_0=(a+b)/2$ and $y_0=(f+g)/2$, then we have
\bean
e^{\frac{c_h(a-x_0)}{2d}}=e^{\frac{c_h(a-b)}{4d}}, \quad e^{\frac{c_h(b-x_0)}{2d}}=e^{\frac{c_h(b-a)}{4d}},
\eean
and 
\bean
{C}_{(n, m)}&=&\sin\left[\frac{n\pi(x_0-a)}{b-a}\right]\sin\left[\frac{m\pi(y_0-f)}{g-f}\right]=\sin\left[\frac{n\pi}{2}\right]\sin\left[\frac{m\pi}{2}\right]\\
&=&(-1)^{n+1}{C}_{(n, m)}.
\eean
Thus the term $(-1)^{n+1}$ does not contribute to the sum of $J_B(t)$ when $n$ is even. Hence, this yields
\[\frac{J_B(t)}{J_A(t)}=e^{\frac{c_h}{2d}(b-a)}.\]
Following this, we immediately have
\bean
N_{B}(t)=\int_0^t J_{B}(t)dt=e^{\frac{c_h}{2d}(b-a)}\int_0^t J_{A}(t)dt=e^{\frac{c_h}{2d}(b-a)}N_A(t),\\
T_{B}(t)=\int_0^{t} t\;J_B(t)dt =e^{\frac{c_h}{2d}(b-a)}\int_0^{t} t\;J_A(t)dt=e^{\frac{c_h}{2d}(b-a)}T_A(t).
\eean
Therefore, the proof is complete.
\end{proof}


Next we turn to the case for which there is no absorbing boundary condition imposed. More specifically, we consider the following diffusion model
\bea\label{two-D-model}
\frac{\partial u}{\partial t}=d\left(\frac{\partial^2u}{\partial x^2}+\frac{\partial^2u}{\partial y^2}\right)-c_h\frac{\partial u}{\partial x}-c_v\frac{\partial u}{\partial y}+\gamma(t) u, \quad (x, y)\in \Real^2,
\eea
with initial condition $u|_{t=0}=\delta(P-P_0)$. In this case, its solution is given by
\bean
u(P, t|P_0)=\frac{1}{4\pi d t}\;e^{-\frac{(x-x_0-c_ht)^2+(y-y_0-c_vt)^2}{4dt}-\int_0^t\gamma(t)dt},\quad t>0
\eean
where $P_0=(x_0, y_0)$.

Since now the absorbing boundary condition is no longer imposed, this setting includes the situations in which  
the individuals can stay on the boundary, or turn back once a boundary has been crossed.

For this case, the population density function is available on each side, thus it allows us to calculate the
population density function on each side directly. For example, on side $A$, its population density at time $t$ can be obtained by
\bean
{\cal N}_A(t)=\int_f^g u(P, t|P_0)\Big|_{x=a} dy 
\eean 
and other sides are computed in the similar way. The mean occupancy time on side $A$ up to $t$ can be defined as
\[{\cal T}_A(t)=\int_0^t{\cal N}_A(t)dt,\]
 and other sides can be defined similarly. For the free boundary case,  the mean first passage time will be challenging to obtain since an insect may crosses a boundary more than once. 
\begin{theorem}\label{thm2} Let $\Omega=\left\{(x, y)|\; a\leq x\leq b, f\leq y\leq g\;\right\}\subset \Real^2$. We denote the side $x=a, x=b$ and $y=f, y=g$ of $\Omega$ by $A$, $B$ and $F$, $G$, respectively. If the initial released point is at the center of $\Omega$, i.e., $x_0=(a+b)/2$ and $y_0=(f+g)/2$, then for $t>0$ we have
\bean
\frac{u(P; t|P_0)|_{x=b}}{u(P; t|P_0)|_{x=a}}=\frac{{\cal N}_B(t)}{{\cal N}_A(t)}=\frac{{\cal T}_B(t)}{{\cal T}_A(t)}=e^{\frac{c_h}{2d}(b-a)},
\eean
and
\bean
\frac{u(P; t|P_0)|_{y=g}}{u(P; t|P_0)|_{y=f}}=\frac{{\cal N}_G(t)}{{\cal N}_F(t)}=\frac{{\cal T}_G(t)}{{\cal T}_F(t)}=e^{\frac{c_v}{2d}(g-f)}.
\eean
\end{theorem}
\begin{proof} Notice that if $x_0=(a+b)/2$ and $y_0=(f+g)/2$, then for $t>0$ we have
\bean
e^{-\frac{(b-x_0-c_ht)^2}{4dt}}= e^{\frac{c_h}{2d}(b-a)}\cdot e^{-\frac{(a-x_0-c_ht)^2}{4dt}},
\eean
which leads to
\bean
u(P; t|P_0)|_{x=b}=e^{\frac{c_h}{2d}(b-a)}\;u(P; t|P_0)|_{x=a}.
\eean
Therefore, we have
\bean
{\cal N}_B(t)=\int_f^g u(P, t|P_0)\Big|_{x=b} dy =e^{\frac{c_h}{2d}(b-a)}\int_f^g u(P, t|P_0)\Big|_{x=a} dy=e^{\frac{c_h}{2d}(b-a)}{\cal N}_A(t),
\eean
and
\bean
{\cal T}_B(t)=\int_0^t{\cal N}_B(t)dt=e^{\frac{c_h}{2d}(b-a)}\int_0^t{\cal N}_A(t)dt=e^{\frac{c_h}{2d}(b-a)}{\cal T}_A(t).
\eean
For the side $F$ and $G$, the proof is similar.
\end{proof}

Therefore, based on Theorem \ref{thm1} and Theorem \ref{thm2}, insect drift movement in a two-dimensional domain is completely governed by the ratios $\frac{c_h}{d}$ and $\frac{c_v}{d}$, and the norm of actual drift is characterized by $\sqrt{c_h^2+c_v^2}/d=|\vec{v}|/d$.

For the free boundary case, Fig. \ref{fig-u-f}, Fig. \ref{fig-N-f}, and Fig. \ref{fig-T-f} show the ratio characteristics similar to the absorbing boundary case.  

\begin{figure}[!htbp]
    \centering
    \includegraphics[width=3in]{Plot_of_u_on_boundary_A_vs_y_t_Xiao_unit_d_01.eps}
    \includegraphics[width=3in]{Plot_of_u_on_boundary_B_vs_y_t_Xiao_unit_d_01.eps}
    \includegraphics[width=3in]{Plot_of_u_on_boundary_C_vs_x_t_Xiao_unit_d_01.eps}
    \includegraphics[width=3in]{Plot_of_u_on_boundary_D_vs_x_t_Xiao_unit_d_01.eps}
    \caption{The plots of $({u}_A(y,t), {u}_B(y,t), {u}_F(x,t), {u}_G(x,t))$ with $d=0.1, c_{h}=0.2, c_{v}=-0.3, N_{0}=1$ on unit square.}
    \label{fig-u-f}
\end{figure}





\begin{figure}[!htbp]
    \centering
    \includegraphics[width=3in]{Plot_Num_of_All_Free_Boundaries_vs_T_unit_d_01_Xiao.eps}
    \includegraphics[width=3in]{Plot_Num_of_All_Free_Boundaries_vs_T_unit_log_d_01_Xiao.eps}
    \caption{The plot of $({\cal N}_A(t), {\cal N}_B(t), {\cal N}_F(t), {\cal N}_G(t))$ with $d=0.1, c_{h}=0.2, c_{v}=-0.3$ on unit square.}
    \label{fig-N-f}
\end{figure}


\begin{figure}[!htbp]
    \centering
    \includegraphics[width=3in]{Plot_MOT_of_All_Free_Boundaries_vs_T_unit_d_01_Xiao.eps}
    \includegraphics[width=3in]{Plot_MOT_of_All_Free_Boundaries_vs_T_unit_log_d_01_Xiao.eps}
    \caption{The plot of $({\cal T}_A(t), {\cal T}_B(t), {\cal T}_F(t), {\cal  T}_G(t))$ with $d=0.1, c_{h}=0.2, c_{v}=-0.3$ on unit square.}
    \label{fig-T-f}
\end{figure}

\clearpage
\section{Concluding Remarks}
In this paper, we present some important characteristics for the insect movement with bias in a rectangular domain. These findings lead to a significant simplification for the estimation of parameters in diffusion models of insect movement.  Statistical experiments show that the estimation based on these new properties is efficient and unbiased.

Although our approach assumes the involved parameters are constant in space, our method is still applicable if these parameters vary but we use a sufficiently small rectangle.  In fact, our methods could be used to map spatial variation in the advection and diffusion parameters, and also determine how these parameters vary with respect to the density of insects, resources, and natural enemies.  This would aid the development of a landscape-level model of insect movement.  

How can these methods be applied to field studies of insect dispersal?  We first suppose that the insects are released at the center of a rectangle.  For most systems, the easiest quantities to measure would likely be the accumulated population across each boundary, at several points in time ($N_A(t_i), N_B(t_i)$, etc.).  These quantities could be measured using sticky panels for flying insects, or pitfall traps for walking ones.  It is not necessary to completely enclose the rectangle (see Section 2.2).  An alternative to a central release would be to initially locate insects in the landscape, and then observe their positions at several points in time within a rectangle marked in some fashion.  If enough insects were observed near this particular initial position, then these observations could be combined to generate values of $N_A(t_i)$ and so forth.  

A decision must also be made as to whether the boundary of the rectangle is absorbing or free, because this alters the method of estimation.  If all insects are trapped as they leave the rectangle, this would obviously be absorbing.  If only some are trapped then the boundary could still be treated as absorbing, however, as long as individuals returning to the rectangle are excluded from the data.  There is some advantage to this procedure because the absorbing boundary case performed better in our simulations.  

If these studies were repeated at multiple locations in space, say at various distances from an attractive resource or habitat patch, it would then be possible to map the spatial variation in the convection and diffusion rates.  This would help determine the actual functional forms of these quantities with respect to distance, a service not provided by existing methods.  One drawback of this procedure is that observations of dispersal need to be made in multiple locations in the landscape.  However, the observations should be less labor intensive than typical for mark-recapture studies, because they are limited in space to the rectangle.  

{\bf We need to update the references.}
\begin{thebibliography}{99}

\bibitem{} Banks, H. T., Kareiva, P. M. and Zia, L. (1988) Analyzing field studies of insect dispersal using two-dimensional transport equations. Environmental Entomology 17: 815-820.  

\bibitem{} Bommarco, R. and Fagan, W. F. (2002) Influence of crop edges on movement of generalist predators: a diffusion approach.  Agricultural and Forest Entomology 4: 21-30.

\bibitem{} Byers, J. A. and Zhang, Q. (2011) Chemical ecology of bark beetles in regard to search and selection of host trees. Pp. 150-190 in Liu, T. and Kang, L. (eds) Recent Advances in Entomological Research.  Higher Education Press, Beijing and Spring-Verlag, Berlin Heidelberg.

\bibitem{} Colazza, S., Cusumano, A., Lo Giudice, D. and Peri, E. (2014) Chemo-orientation responses in hymenopteran parasitoids induced by substrate-borne semiochemicals.  BioControl 59: 1-17.

\bibitem{} Fagan, W. F. (1997) Introducing a ``boundary-flux'' approach to quantifying insect diffusion rates. Ecology 78: 579-587

\bibitem{} Kareiva, P. (1982) Experimental and mathematical analyses of herbivore movement: quantifying the influence of plant spacing and quality on foraging discrimination. Ecological Monographs 52: 261-282.

\bibitem{} Kareiva, P. and Odell, G. (1987) Swarms of prey exhibit ``preytaxis'' if individual predators use area-restricted search.  American Naturalist 130: 233-270.

\bibitem{} Ovaskainen, O. (2004) Habitat-specific movement parameters estimated using mark-recapture data and a diffusion model.  Ecology 85: 242-257.

\bibitem{} Ovaskainen, O., Rekola, H., Meyke, E. and Arjas, E. (2008) Bayesian methods for analyzing movements in heterogeneous landscapes from mark-recapture data.  Ecology 89: 542-554.

\bibitem{Pi98} Pinsky, MA (1998) Partial differential equations and boundary-value problems with applications, American Mathematical Society.

\bibitem{} Reeve, J. D. and Cronin, J. T. (2010) Edge behaviour in a minute parasitic wasp. Journal of Animal Ecology 79: 483-490.

\bibitem{} Reeve, J. D., Cronin, J. T. and Haynes, K. J. (2008) Diffusion models for animals in complex landscapes: incorporating heterogeneity among substrates, individuals and edge behaviours.  Journal of Animal Ecology 77: 898-904.

\bibitem{} Ries, L. and Sisk, T. D. (2008) Butterfly edge effects are predicted by a simple model in a complex landscape. Oecologia 156: 75-86.  

\bibitem{} Schultz, C. B. and Crone, E. E. (2001) Edge-mediated dispersal behavior in a prairie butterfly.  Ecology 82: 1879-1892.

\bibitem{} Turchin, P. and Simmons, G. (1997) Movements of animals in congregation: an Eulerian analysis of bark beetle swarming.  Pp. 113-125 in Parrish, J. K., Hamner, W. M. and Prewitt, C. T. (eds). Animal Groups in Three Dimensions.  Cambridge University Press, New York, NY.

\bibitem{XRXC} M.~Xiao, J.~D.~Reeve, D.~Xu, and J.~T.~Cronin (2013) {\em Estimation of the diffusion rate and crossing probability for biased edge movement between two different types of habitat}, Journal of Mathematical Biology: Volume 67, Issue 3, pp.~ 535-567.

\end{thebibliography}
\end{document}

